% SPDX-License-Identifier: Apache-2.0
%
% A reservation station in TxHDL: the tutorial on the part. Built by
% //docs:station. Every listing is included from the tree and every
% printout and figure is what the build made by running the example.
\documentclass[journal,9pt]{IEEEtran}

% Listings of Rust set by rustfmt at 80 columns: 5.6pt is what fits
% the frame of a column (issue 1061).
\newcommand{\listingsize}{\fontsize{5.6}{6.6}\selectfont}
% SPDX-License-Identifier: Apache-2.0
%
% The house preamble, shared by every document in this directory. Loaded
% after \documentclass, before \title. Keeping it in one file is what
% stops two articles drifting apart in font, listing style or figure
% style.
% Computer Modern. IEEEtran selects Times (ptm) by default, so the face has
% to be chosen explicitly.
%
% Latin Modern is the Computer Modern variety used here, and the choice is
% forced rather than aesthetic. Under T1 encoding, plain `cmr` has no Type 1
% outlines unless cm-super is installed, and it is not in the pinned
% distribution, so pdflatex falls back to EC bitmaps and every glyph in the
% PDF becomes a Type 3 font. That was measured: the first build produced a
% document with 10 Type 3 fonts and no Type 1. Latin Modern is Computer
% Modern redrawn with a full T1 repertoire and real .pfb outlines, so the
% same design comes out scalable.
\usepackage[T1]{fontenc}
\usepackage[utf8]{inputenc}
\usepackage{lmodern}
% microtype is not in the pinned TeX distribution. emergencystretch alone
% keeps the two-column measure from producing overfull lines.
\setlength{\emergencystretch}{3em}

\usepackage{amsmath}
\usepackage{amssymb}
\usepackage{amsthm}
\usepackage{booktabs}
\usepackage{array}
% enumitem is not in the pinned TeX distribution; the list spacing below
% does what \setlist would have done.
\usepackage{float}
\usepackage{xcolor}
\usepackage{listings}
\usepackage{tikz}
\usetikzlibrary{arrows.meta,positioning,fit,backgrounds,calc}
\usepackage[hidelinks,bookmarks=true]{hyperref}

% Numbered, definition-style box for a rule the reader is meant to apply.
\theoremstyle{definition}
\newtheorem{rulebox}{Rule}
\newtheorem{example}{Example}

% Unnumbered asides.
\theoremstyle{remark}
\newtheorem*{tip}{Tip}
\newtheorem*{pitfall}{Pitfall}

% Tighter lists, in place of enumitem's \setlist.
\setlength{\topsep}{2pt}
\setlength{\itemsep}{1pt}
\setlength{\parsep}{0pt}

% Code in the running text. The typewriter face under T1 ligates `<<`
% and `>>` into guillemets, so a shift printed as \code{<<} came out as
% a French quotation mark (issue 571). microtype's \DisableLigatures is
% not in the pinned distribution, but the primitive it rests on is
% pdfTeX's own: \pdfnoligatures strips every ligature from the font it
% is given, and the current font inside \texttt is the typewriter face
% at the size in use, so each size is stripped the first time \code
% sets it. Code wants no ligature at all: two quotes are two quotes.
\newcommand{\code}[1]{\texttt{\ifdefined\pdfnoligatures\pdfnoligatures\font\fi #1}}
\newcommand{\term}[1]{\emph{#1}}

% Syntax highlighting. Four roles, distinguishable in colour and still
% distinguishable in greyscale, because an IEEE article gets printed: the
% keyword is the only bold role, the comment is the only italic one, and the
% two remaining roles differ in lightness as well as in hue.
\definecolor{kwblue}{RGB}{20,60,150}
\definecolor{cmtgreen}{RGB}{90,120,90}
\definecolor{strred}{RGB}{150,50,40}
\definecolor{numgray}{gray}{0.45}
\definecolor{framegray}{gray}{0.70}
\definecolor{bgshade}{gray}{0.97}
% A darker fill for figure cells. The listing background has to stay very
% light to keep code readable; a figure cell has to be clearly shaded or the
% distinction it draws is invisible in print.
\definecolor{cellshade}{gray}{0.90}

% The listing size is the document's choice, because it depends on the
% column: a two-column article sets `\listingsize` to 6pt and keeps its
% listings within 70 columns, or to 5.6pt when it includes source that
% rustfmt sets at 80, which is what fits a column's frame (issue 1061);
% a one-column document keeps the body size and includes source files
% at 80. `//docs:frame_test` checks every listing line against its frame. Lines are never broken; a line that does not fit
% overflows the frame, which is visible, rather than wrapping, which is
% not.
\providecommand{\listingsize}{\scriptsize}

% `'` is deliberately not a string delimiter: the language uses it for loop
% labels, as Rust does, so treating it as a quote would swallow the rest of
% the line.
\lstdefinestyle{txhdl}{
  basicstyle=\ttfamily\listingsize,
  keywordstyle=\bfseries\color{kwblue},
  commentstyle=\itshape\color{cmtgreen},
  stringstyle=\color{strred},
  numberstyle=\tiny\color{numgray},
  morekeywords={mod,use,pub,const,type,struct,enum,trait,impl,fn,async,let,mut,
    match,if,else,loop,while,for,in,break,continue,return,as,await,
    module,bus,channel,transaction,protocol,pipeline,flow,seq,config,
    port,stage,pipe,instance,connect,bind,tag,domain,blackbox,foreign,
    property,role,signal,handler,true,false,
    sequence,interface,modport,implements,package,import,var,wire,func,proc,
    case,default,next,fork,branch,join,state,fsm},
  morecomment=[l]{//},
  morestring=[b]",
  showstringspaces=false,
  columns=fullflexible,
  keepspaces=true,
  breaklines=false,
  % Framed, so a listing is bounded rather than floating in the column.
  frame=single,
  rulecolor=\color{framegray},
  framesep=3pt,
  backgroundcolor=\color{bgshade},
  xleftmargin=3pt,
  xrightmargin=3pt,
  aboveskip=6pt,
  belowskip=6pt,
  % Every listing is numbered and captioned, so it can be referred to by
  % number. `caption` without `float` numbers the listing and leaves it
  % where it was written, which is what a listing in running prose needs.
  captionpos=b,
  abovecaptionskip=2pt,
  belowcaptionskip=6pt,
}

% Figures drawn as text. Framed the same way, and with the highlighting
% switched off, because a box-and-arrow diagram is not code and half its
% words turning blue reads as an accident. Setting a style adds keys rather
% than clearing the ones already set, so the three roles are neutralised by
% hand rather than by leaving them out. Program output is printed in this
% style too, and a netlist's assignment can run past the frame, so a line
% that does not fit is broken and the rest indented; source never is.
\lstdefinestyle{ascii}{
  basicstyle=\ttfamily\listingsize,
  keywordstyle=\ttfamily,
  commentstyle=\ttfamily,
  stringstyle=\ttfamily,
  columns=fullflexible,
  keepspaces=true,
  breaklines=true,
  breakindent=2em,
  frame=single,
  rulecolor=\color{framegray},
  framesep=3pt,
  backgroundcolor=\color{bgshade},
  xleftmargin=3pt,
  xrightmargin=3pt,
  aboveskip=6pt,
  belowskip=6pt,
}
\lstset{style=txhdl}

% House TikZ styles. `shorten` lives on the arrow styles rather than on each
% arrow, so every arrow in the document gets the gap, including ones added
% later. TikZ clips a path at a node's border, so without it an arrowhead
% butts against the box with no gap at all. Keep `node distance` well above
% twice the shorten, or the arrow shrinks to nothing.
\tikzset{
  box/.style={draw,rounded corners=2pt,align=center,inner sep=4pt,
              font=\footnotesize},
  boxf/.style={box,fill=bgshade},
  lbl/.style={font=\scriptsize\itshape,align=center},
  fl/.style={-{Stealth[length=4pt]},shorten >=2.5pt,shorten <=2.5pt},
  fld/.style={fl,dashed},
}


% The contents list: the class gives a section's number 2.75em, which
% a Roman numeral past thirty overruns onto the title, so the box is
% wider here, and a subsection's entry starts where a title does.
\makeatletter
\def\l@section#1#2{\addpenalty{\@secpenalty}\addvspace{1.0em plus 1pt}%
    \@tempdima 5.75em \begingroup \parindent \z@ \rightskip \@pnumwidth%
    \parfillskip-\@pnumwidth {\bfseries\leavevmode #1}\hfil\hbox to\@pnumwidth{\hss #2}\par%
    \endgroup}
\def\l@subsection{\@dottedtocline{2}{5.75em}{5em}}
\def\l@subsubsection{\@dottedtocline{3}{10.75em}{5.5em}}
\makeatother


\InputIfFileExists{buildstamp.tex}{}{}
\providecommand{\buildstamp}{Draft build (outside the Bazel flow).}

\title{A Reservation Station in TxHDL}

\author{Filip Filmar and Dragiša Janković%
\thanks{Both authors are independent. Filip Filmar is the
corresponding author, at f@hdlfactory.com; Dragiša Janković is
reached at linkedin.com/in/dragisa-jankovic.}%
\thanks{This tutorial demonstrates the reservation station component from the
parts library, for developers who have completed the introductory
tutorial~\cite{tutorial} and seek to utilize or study out-of-order execution
primitives. Listings are extracted directly from the repository \code{HDL/txhdl}
at build time; logs and waveforms reflect actual test execution. The component
source appears in~\cite{runtime}, reference examples in~\cite{examples}, and the
project index~\cite{cover} catalogues all companion documents.}%
\thanks{The authors used a large language model, Claude, as an
assistant in exploring concepts and in drafting and constructing
the documents and implementations; every repository commit documents
this collaboration and includes the prompt sequence verbatim.}%
\thanks{\buildstamp}}

\markboth{A Reservation Station in TxHDL}{Filmar and Janković: A Reservation Station in TxHDL}

\begin{document}
\maketitle

\begin{abstract}
A reservation station gathers asynchronous operands for an operation arriving
out of order across independent channels, dispatching them downstream once all
dependencies resolve. The TxHDL parts library provides parameterised
reservation stations, from \code{Station2} to \code{Station10}. Each variant
lowers to synthesizable RTL exposing one input channel per operand and a unified
output channel. This document describes the hardware mechanics of the station,
traces its operational rules cycle by cycle against an executable design,
details FIFO queue integration, and analyzes the emitted netlist.
\end{abstract}

\begin{IEEEkeywords}
Reservation station, Tomasulo algorithm, out-of-order execution, queues,
hardware description languages, TxHDL.
\end{IEEEkeywords}

\section{What It Is}
\label{sec:s-what}

Operations requiring multiple operands must synchronize awaiting the final
input. When operands originate from independent producers across distinct
channels in non-deterministic order, intermediate storage must buffer early
arrivals while isolating distinct operations. Tomasulo reservation stations
solve this coordination problem. Each operation is assigned a unique tag. Each
incoming operand holds the tag of its parent operation. The station allocates a
storage line per tag, subdivided into cells for each input channel. Operands
latch into their assigned line cells. When all cells within a line hold valid
operands, the line completes and dispatches all values concurrently alongside
its tag.

In TxHDL, input and output interfaces operate as transactional channels. An
input asserts \code{ready} when the station has buffer capacity for the offered
transaction, determined by whether the targeted tag cell is unoccupied. Output
\code{valid} indicates a completed line ready for dispatch, while output
\code{ready} reflects downstream buffer availability.

\section{The Interface}
\label{sec:s-interface}

Input channels transport \code{Tagged<TB, T>} structures, packaging a
\code{TB}-bit tag alongside payload data of type \code{T}. This struct derives
both \code{Transaction} and \code{Value}. An $N$-input reservation station is
instantiated as \code{StationN<TB, L, T0, .., TN>}, parameterised by tag width
\code{TB}, line capacity \code{L} (equal to \code{1 << TB}), and individual
operand types \code{T0} through \code{TN}. The \code{run} method accepts a tuple
of $N$ input receivers, \code{Rx<Tagged<TB, Ti>>}, and a single output
transmitter, \code{Tx<LineN<TB, T0, .., TN>>}. The emitted \code{LineN} struct
bundles the tag with resolved operand values \code{v0} through \code{vN}.
Listing~\ref{lst:s-part} presents the library source code: the tagged payload
definition and nine concrete station generators produced by macro
\code{station!}.

\lstinputlisting[rangebeginprefix=//\ begin\{,rangebeginsuffix=\},rangeendprefix=//\ end\{,rangeendsuffix=\},includerangemarker=false,linerange=part-part,caption={From \code{lib/parts/src/station.rs}: the tagged payload and the nine stations; the rule is in the module's own words there, and its tests follow.},label={lst:s-part}]{lib/parts/src/station.rs}

\section{The Rule}
\label{sec:s-rule}

In each clock cycle, the station examines incoming channel offers and evaluates
its internal state: a dedicated memory array storing operand cells per input,
and an occupancy mask tracking valid cells per line. The dispatch discipline
operates as follows:

\begin{enumerate}
\item A cell is \emph{free} when its corresponding bit in the input occupancy
  mask is clear.
\item Consuming an input \emph{completes} a line when all other operand cells in
  that line already hold valid data, or when another input offers the matching
  tag concurrently into a free cell.
\item The \emph{sent} line corresponds to the lowest-indexed completing input,
  provided the downstream channel asserts \code{ready}. At most one line
  dispatches per cycle.
\item An input is \emph{taken} if its cell is free and consumption either
  completes no line or completes the active line being dispatched. Inputs whose
  acceptance would complete a line that cannot dispatch this cycle must wait,
  occurring when downstream channels lack capacity or when a lower-indexed
  input claims dispatch priority.
\item Taken inputs write their payload into cell memory and set their occupancy
  bit. Dispatched lines clear their occupancy bit across all masks; output
  values are routed from memory cells, with concurrent inputs forwarded
  directly through bypass logic.
\end{enumerate}

Under this discipline, inputs never block on cells belonging to other inputs;
they wait only on their own cell capacity or downstream backpressure. If an input
repeatedly presents tags whose matching lines never complete due to stalled
upstream producers, the channel stalls indefinitely. The station provides no
internal timeout or eviction arbitration.

\section{The Rule, Cycle by Cycle}
\label{sec:s-cycles}

The reference example configures a \code{Station2} module coordinating an 8-bit
byte and a 16-bit halfword with a 2-bit tag, supporting four lines. The first
source generates bytes tagged 0, 1, 2, 3, 0, 2 across successive cycles as the
station permits. The second source offers halfwords tagged 2, 0, 3, 1, 2, 0,
pausing every third cycle. The receiving sink accepts output lines on two
cycles out of four. Listing~\ref{lst:s-out} records the cycle-by-cycle execution
log, documenting offered inputs, occupancy masks, and retiring lines. Offers
arrive at station inputs in the cycle following transmission, reflecting
standard register-transfer delays.

\lstinputlisting[style=ascii,caption={What \code{ex\_station} printed when the build ran it: the offers, the masks, the lines taken; then the lowering.},label={lst:s-out}]{out_station.txt}

Trace execution across successive cycles:

\begin{itemize}
\item Cycle 1: Source 0 offers byte tag 0; source 1 offers halfword tag 2.
\item Cycle 2: The station consumes both into cell 0 of mask 0 and cell 2 of
  mask 1 (masks: \code{0001} and \code{0100}).
\item Cycle 3: Halfword tag 0 completes line 0 because byte cell 0 is valid.
  Line 0 dispatches immediately, while byte tag 1 fills cell 1 (masks:
  \code{0010} and \code{0000}). The sink consumes line 0 in cycle 4.
\item Cycle 4: Byte tag 2 completes line 2 because halfword cell 2 is valid.
  Line 2 dispatches and cell 2 clears (masks: \code{0000} and \code{0000}).
\item Cycle 5: Both sources offer tag 3 simultaneously. Both complete line 3
  concurrently, dispatching immediately without updating storage masks.
\item Cycle 6: Byte tag 0 fills cell 0 (\code{0011}). Halfword tag 1 would
  complete line 1, but downstream output buffers retain pending lines 2 and 3;
  lacking capacity, halfword 1 waits.
\item Cycle 7: Halfword 1 waits again; byte tag 2 fills cell 2 (\code{0111}).
\item Cycle 8: The sink consumes line 2; output space opens, line 1 dispatches,
  and cell 1 clears (\code{0101}).
\item Cycle 9: Halfword tag 2 completes line 2 (\code{0001}).
\item Cycle 11: Output backpressure stalls halfword tag 0 until cycle 12,
  completing line 0 once masks empty. The sink consumes lines 1, 2, and 0 in
  cycles 11, 12, and 15.
\end{itemize}

Figure~\ref{fig:s-wave} illustrates the complete waveform: input valid/ready
handshakes, mask transitions, and output retirements. Input 1 asserts backpressure
through cycles 6, 7, and 11, matching the waiting conditions of Rule 4.

\begin{figure*}[t]
\centering
\input{station_timing.tex}
\caption{The example's run: the two inputs' tags, values and
handshakes, the masks, and the lines out with their values.}
\label{fig:s-wave}
\end{figure*}

\section{Using a Station}
\label{sec:s-use}

Listing~\ref{lst:s-example} presents the complete top-level integration. The
module instantiates \code{Station2<2, 4, U<8>, U<16>>} directly. Input ports
attach to two channel receivers, and output ports connect to a single
transmitter. Source processes emit \code{Tagged} structures while sink processes
consume \code{Line2} values, accessing fields \code{tag}, \code{v0}, and
\code{v1}. Higher-order stations extend this structure up to \code{Station10}.
Tags require at least two bits; single-bit values map to scalar wires in VHDL
and cannot index multi-line memories.

\lstinputlisting[caption={\code{ex\_station.rs}. Run with \code{bazel run //lib/examples:ex\_station}.},label={lst:s-example}]{lib/examples/ex_station.rs}

Because native channels buffer two items, the station output tolerates two
cycles of downstream backpressure before stalling. When consumers process
traffic in large bursts, a dedicated \code{Fifo} queue is placed downstream of
the station: \code{Fifo<Line2<2, U<8>, U<16>>, 3, 8>} provides an 8-entry queue
indexed by 3 address bits. Listing~\ref{lst:s-queue} demonstrates FIFO operation
buffering nine incoming items across a 4-entry queue.

\lstinputlisting[style=ascii,caption={What \code{ex\_queue} printed when the build ran it: a FIFO of four between two channels, filling and draining; then its lowering.},label={lst:s-queue}]{out_queue.txt}

\section{The Same Station in Verilog}
\label{sec:s-verilog}

Listing~\ref{lst:s-v} provides a hand-written Verilog implementation of the
two-input station for comparison. Both modules implement identical port
signatures, with one distinction: TxHDL netlists include an implicit \code{rst}
input. The Verilog source represents an independent translation of the
operational rules in Section~\ref{sec:s-rule}. In the testbench, both units
execute in parallel under identical input stimulus, verifying cycle-by-cycle
output equivalence. Listing~\ref{lst:s-src} displays macro-expanded Rust source
from \code{station!}, and Listing~\ref{lst:s-sizes} compares physical line
counts.

\lstinputlisting[style=ascii,caption={\code{station.v}: the two-input station by hand.},label={lst:s-v}]{lib/examples/station.v}

\lstinputlisting[caption={\code{station2.rs}: the two-input station as \code{station!} wrote it, printed by \code{ex\_station --source}.},label={lst:s-src}]{station2.rs}

\lstinputlisting[style=ascii,caption={The two, by lines, as the build counted them.},label={lst:s-sizes}]{station_sizes.txt}

Both designs implement identical hardware behavior with matching net names.
Key differences emerge in authoring overhead and verification guarantees:

\begin{itemize}
\item \textbf{Bit-Widths and Slicing:} Verilog requires explicit bit ranges
  (\code{[9:0]}, \code{[17:0]}) and manual payload extraction
  (\code{in0\_data[9:8]}); altering a data type mandates coordinated edits
  across five locations. TxHDL extracts fields automatically from typed structs
  (\code{head0.tag}), parameterising types cleanly. While Verilog supports
  parameterised widths, it cannot parameterise port counts cleanly; the TxHDL
  macro generates all nine variants from one template.
\item \textbf{Handshake Protocols:} Verilog authors manually coordinate
  \code{ready} and \code{valid} nets, avoiding combinational feedback loops that
  cause synthesis errors. TxHDL authors specify channel operations via
  \code{recv\_if(take0)} and conditional \code{send}, allowing lowering passes
  to derive valid/ready handshakes automatically.
\item \textbf{Verification Coverage:} Verilog requires handwritten testbenches.
  The TxHDL station executed as pure software for 400 cycles before lowering,
  and its emitted netlists are verified continuously under both \code{nvc} and
  Verilator on every build.
\item \textbf{Physical Line Counts:} Verilog appears shorter because comments
  are omitted in Listing~\ref{lst:s-sizes}. However, the Verilog author writes
  all 53 lines manually for every port configuration. The TxHDL designer
  invokes \code{station!(Station2, 2)} once, reusing the definition across
  arbitrary widths and types.
\end{itemize}

\section{The Netlist}
\label{sec:s-netlist}

Listing~\ref{lst:s-out} concludes with the synthesized Verilog netlist for the
two-input station. Each mask lowers to a 4-bit register, and operand cells map
to 4-word synchronous memories. Intermediate signals directly mirror rule
terminology: \code{cell\_free0}, \code{completes0}, \code{line\_tag},
\code{send\_line}, and \code{take0}. Tags and payloads extract from input vectors
at struct-defined offsets; one-hot masks calculate via bit shifts; and output
data concatenates the tag with resolved operand values. The lowering engine
produces equivalent VHDL entities using \code{shift\_left} and \code{to\_integer}.
Both netlists are co-simulated against the Rust execution trace on every build.

\begin{thebibliography}{9}

\bibitem{tutorial}
F.~Filmar and D.~Janković, ``TxHDL, step by step,'' project
repository \code{HDL/txhdl}, \code{//docs:tutorial}, 2026.

\bibitem{runtime}
F.~Filmar and D.~Janković, ``The TxHDL runtime,'' project repository
\code{HDL/txhdl}, \code{//docs:runtime}, 2026.

\bibitem{examples}
F.~Filmar and D.~Janković, ``TxHDL by example,'' project repository
\code{HDL/txhdl}, \code{//docs:examples}, 2026.

\bibitem{cover}
F.~Filmar and D.~Janković, ``TxHDL documents,'' project repository
\code{HDL/txhdl}, \code{//docs:cover}, 2026.

\end{thebibliography}

\end{document}
